\begin{thebibliography}{9}
% The {9} works for 1 to 9 items to cite.
% If you have 10 to 99 items, use {99}.
% If you have 100 to 999 items, use {999}.
% If you have more than 1000 items, you need to get a new 
% hobby and/or a significant other.

\bibitem{LA1979}
	Lars Ahlfors,
	\emph{Complex Analysis},
	McGraw-Hill, New York, 1979.
\bibitem{FL1882}
	Ferdinand von Lindemann,
	\"{U}ber die Zahl $\pi$,
	\emph{Mathematische Annalen} \textbf{20} (1882), 213--225.
\bibitem{AW1995}
	Andrew Wiles,
	Modular Elliptic Curves and Fermat's Last Theorem,
	\emph{Ann. Math.} \textbf{141} (1995), 443-551.
\end{thebibliography}

% The convention here is to use for each label the author's initials
% and the year of publication, but you may feel free to use whatever
% is easiest for you to remember.


